%! Simplifying Rational Expressions — Unit 5, Lesson 5.1 — Homework (Answer Key)
\documentclass[10pt]{article}
\usepackage{algebra2-article}
\usepackage{algebra2-key}   % pulls in -boxes; do NOT also load -boxes
\newcommand{\TallMath}[1]{$\displaystyle #1\rule[-1.4em]{0pt}{3.2em}$}
\newcommand{\lgrid}[4]{%
  \draw[step=1, linegray!40, very thin] (#1,#3) grid (#2,#4);
  \draw[->, semithick, charcoal] (#1,0)--(#2,0) node[below right, font=\scriptsize]{$x$};
  \draw[->, semithick, charcoal] (0,#3)--(0,#4) node[left, font=\scriptsize]{$y$};}

\begin{document}

\pageheader{Unit 5, Lesson 5.1}{Homework --- Answer Key}
\namedateperiod

\vspace{0.10in}

\begin{notesbox}{Practice}
Every answer needs \emph{two} parts: the simplest form \textbf{and} the restrictions, taken from the
\textbf{original} denominator.

\begin{enumerate}[leftmargin=1.9em, itemsep=11pt]
  \item \textbf{Monomial factors.} Simplify and restrict.
    \begin{enumerate}[label=(\alph*), leftmargin=2.1em, itemsep=5pt]
      \item $\dfrac{14x^5}{21x^3}=\;$\ans{$\frac{2x^2}{3}$} \quad $x\neq$ \ans{0}
      \item $\dfrac{8a^2b^3}{12a^5b}=\;$\ans{$\frac{2b^2}{3a^3}$} \quad restrictions: \ans{$a\neq0,\ b\neq0$}
      \item $\dfrac{5x^2y}{15xy^3}=\;$\ans{$\frac{x}{3y^2}$} \quad restrictions: \ans{$x\neq0,\ y\neq0$}
    \end{enumerate}

  \item \textbf{Binomial and quadratic factors.} Factor completely first. One of these has
        \emph{opposite} binomials --- watch the sign.
    \begin{enumerate}[label=(\alph*), leftmargin=2.1em, itemsep=7pt]
      \item $\dfrac{x^2-49}{x-7}=\;$\ans{$x+7$} \quad $x\neq$ \ans{7}
      \item $\dfrac{5x-20}{x^2-16}=\;$\ans{$\frac{5}{x+4}$} \quad $x\neq$ \ans{$4,\ -4$}
      \item $\dfrac{x^2+7x+10}{x^2+3x-10}=\;$\ans{$\frac{x+2}{x-2}$} \quad $x\neq$ \ans{$2,\ -5$}
      \item $\dfrac{3-x}{x^2-9}=\;$\ans{$-\frac{1}{x+3}$} \quad $x\neq$ \ans{$3,\ -3$}
      \item $\dfrac{2x^2-8x}{x^2-16}=\;$\ans{$\frac{2x}{x+4}$} \quad $x\neq$ \ans{$4,\ -4$}
      \item $\dfrac{x^2-6x+9}{x^2-9}=\;$\ans{$\frac{x-3}{x+3}$} \quad $x\neq$ \ans{$3,\ -3$}
    \end{enumerate}

  \item \textbf{Match the algebra to the picture.} The graph below is $g(x)=\dfrac{x^2+2x-8}{x-2}$.
    \begin{center}
    \begin{tikzpicture}[scale=0.44]
      \lgrid{-7}{4}{-3}{8}
      \begin{scope}
        \clip (-7,-3) rectangle (4,8);
        \draw[very thick, forest] (-7,-3)--(4,8);
      \end{scope}
      \draw[forest, very thick, fill=white] (2,6) circle (4.2pt);
    \end{tikzpicture}
    \end{center}
    Factored form: \ans{$\frac{(x+4)(x-2)}{x-2}$} \quad Simplest form: \ans{$x+4$}, \; for $x\neq$ \ans{2}
    \\[3pt]
    The hole is at \ans{$(2,6)$}. \quad Does $g$ have a vertical asymptote? \ans{No} \\[3pt]
    Why is the graph a \emph{straight line} even though $g$ is a rational function?
    \ansline{the only denominator factor cancels, so away from $x=2$ the expression \emph{is} the line $y=x+4$}

  \item \textbf{Justify.} Are $\dfrac{x^2-1}{x-1}$ and $x+1$ the same expression? Answer in terms of
        \emph{inputs}: where do they agree, where do they disagree, and what is the graph of each?
        \ansline{They are \emph{equivalent} --- $\frac{(x-1)(x+1)}{x-1}=x+1$ --- and they agree at every input except}
        \ansline{$x=1$, where the first is undefined ($\frac00$) and the second equals $2$. So the graph of $x+1$ is the}
        \ansline{whole line, while the graph of $\frac{x^2-1}{x-1}$ is that same line with a hole punched out at $(1,2)$.}
\end{enumerate}
\end{notesbox}

\begin{extensionbox}
\textbf{Part 1 --- Three students, three answers.} Each student simplified
$\dfrac{x^2-4}{x^2+4x+4}$ and got something different. Decide who is right, and name each error.
\begin{center}
\renewcommand{\arraystretch}{1.4}
\begin{tabularx}{\linewidth}{@{}l l X@{}}
\hline
\textbf{Student} & \textbf{Answer} & \textbf{Right or wrong --- and why?} \\
\hline
Amara & $\dfrac{x-2}{x+2}$, $x\neq-2$ & \ans{Correct: $\frac{(x-2)(x+2)}{(x+2)^2}$, one $(x+2)$ divides out.} \\
Ben   & $\dfrac{-4}{4x+4}$ & \ans{Wrong: cancelled the $x^2$ \emph{terms}, which are added, not factors.} \\
Cruz  & $\dfrac{x-2}{x+2}$, no restrictions & \ans{Incomplete: right form, but $x=-2$ is still excluded.} \\
\hline
\end{tabularx}
\end{center}

\vspace{4pt}
\textbf{Part 2 --- Design one (preview of Lesson 5.5).} Write a rational expression whose graph has a
\textbf{hole} at $x=2$ and a \textbf{vertical asymptote} at $x=-1$.
\begin{center}
\ans{$\frac{(x-2)(x+3)}{(x-2)(x+1)}$ \; (any numerator factor other than $x+3$ also works)}
\end{center}
\begin{enumerate}[label=(\alph*), leftmargin=1.9em, itemsep=5pt]
  \item Which factor did you have to put on \emph{both} the top and the bottom, and why?
  \item Which factor appears \emph{only} on the bottom, and why does that make a wall instead of a
        hole?
  \item What is the $y$-coordinate of your hole?
\end{enumerate}
\ansline{(a) $(x-2)$ --- it must be in the denominator to exclude $x=2$, and in the numerator so it cancels,}
\ansline{which is exactly what turns the break into a single missing point. (b) $(x+1)$ stays in the denominator with}
\ansline{nothing to cancel it, so near $x=-1$ the outputs grow without bound --- a vertical asymptote. (c) substitute $x=2$ into the simplified $\frac{x+3}{x+1}$: the hole is at $\left(2,\frac53\right)$.}
\end{extensionbox}

\begin{spiralbox}
\textbf{Coming next --- Lesson 5.2: Multiplying \& Dividing Rational Expressions.} Today's four steps
do not go away; they become step one of everything that follows. To multiply, you factor
\emph{everything} first and then divide out across the whole product; to divide, you keep-change-flip
and then do the same. The one new wrinkle: when you flip a divisor, its \emph{numerator} moves to the
bottom --- so it starts generating restrictions too. Keep tonight's habit of listing restrictions
before you cancel; in 5.2 there will be three denominators to watch instead of one.
\end{spiralbox}

\begin{teachernote}
\textbf{Item 2(d)} is the sign check: $\frac{3-x}{x^2-9}=\frac{-(x-3)}{(x-3)(x+3)}=-\frac{1}{x+3}$.
An answer of $\frac{1}{x+3}$ means the opposite binomials were treated as identical --- the most common
loss of a point on this set. \textbf{Item 2(f)} tests the perfect-square trinomial: only \emph{one}
$(x-3)$ cancels, so $x=3$ is a hole and $x=-3$ is still a wall.

\vspace{3pt}
\textbf{Item 3} is worth spending debrief time on: a rational function whose graph is a
\emph{line} still unsettles students, and it is the cleanest possible picture of what cancelling does.
Ask why the hole is at height $6$ and not at $0$.

\vspace{3pt}
\textbf{Extension Part 2} is a genuine preview of 5.5 (build a graph from an equation) and of 5.6
(extraneous solutions). Accept any numerator factor that keeps $(x-2)$ on both levels; the $y$-value of
the hole then changes with it, which is a good thing to surface --- the hole's location is
\emph{not} determined by the restriction alone.
\end{teachernote}

\end{document}
